\chapter{第四集 \quad 古国}

\section{正文}

蜿蜒清翠的天目山向着衔接于东麓的杭嘉湖平原，伸出两支余脉，如一双臂膀环拥着余杭山水。

5300年前在散落于平原的孤峰残丘之上，村落连绵迭起，人们在山林狩猎采集，在湿地种植水稻。当人口逐渐蕃息，耕作技术发展，村庄和稻田便向平原深处蔓延。为了防备洪涝灾患，先民们堆土成墩，居住在墩上，死后也葬在墩上。取土时形成的低洼地带被顺势整治成河陂，再与自然水网连通便成为来往运输的水道。

时光荏苒，三四百年后这里渐渐成为方圆百里间的中心，良渚古城拔地而起。

4900年前的夏秋之际，一场大雨倾泻而下，袭击了整个天目山区。掌握尖端制玉技术的良渚匠人正在雕刻一件玉琮。神徽的纹理细如发丝，稍有差错便前功尽弃。这时屋外突然电闪雷鸣，眼前变得晦暗不清，他知无法再继续只得搁下手中的刻刀。

此时的莫角山上，王族们隔着雨幕向远处望去，忧心忡忡……连日暴雨威胁着良渚城的安危。在漫长的时光中，水滋润着长江下游这片丰腴的土地，发展出高度发达的稻作农业。洗练出自马家浜始，经岭家滩、松泽至良渚这样数千年间一脉相承的区域文化。然而，水也曾一次次化身为惊涛和洪涝，阻断先民迈向文明的步伐。是以 5000年前良渚先民决定在这世代享溪的水乡平原上修建这座城址时，选址和布局都经过精心的考量。

良渚城三面环山形成天然屏翳，东边的腹地则平畴千里，水网交织，是再好不过的宜居之地。只是夏季多雨将导致北面山洪倾泻；秋天伏旱又会使水稻绝收、河道干涸。

良渚城的营建者们深知要想突破发展的桎梏，必须解决这频发的水患困扰。于是他们开始构画一个宏大的水坝系统的蓝图，经考古发现的11条堤坝遗址构成了具有上下游两级水库的完整水利系统。

西北方向大遮山中的岗公岭、老虎岭、周家畈、秋坞、石坞、蜜蜂垄组成了高水坝群；西南方向狮子山、鲤鱼山、官山、梧桐弄组成了低水坝群；东南方向塘山是平行于大遮山走向的山前长堤。

良渚堤坝的结构类似于现代的黏土心墙坝，坝体外侧用黄色黏土作为坝壳，坝体内则填筑淤泥和草裹泥。"草裹泥"就是用南荻等植物茎秆包裹，再用竹条进行绑扎固定的块状泥土，类似于现代抗洪时使用的编织袋装土。将一块块的"草裹泥"纵横交叠，堆筑成堤，既能提高坝体的强度，也能加快固结过程，在运输和垒筑时更是便于分工协作。

良渚人还巧妙地利用自然环境的特征，对城外山体进行改造，从新发掘的这段石岭头坝体能清楚地看到长堤的迎水面处采用了石砌结构，因此能承受住剧烈水流的冲击。

汛期时山洪被长堤引向西侧，高低二级水坝形成的库区中，可避免良渚城所在的山间平原出现洪灾；伏旱期时库区又可以向城里城外供水，解决农田灌溉和生活用水的需求，使都城的交通命脉保持畅通。这是同时期世界上规模最大的水坝系统，蓄水量达4600余万立方米，相当于3个西湖。

4900年前那场持续多日的大雨终于停息，精密设计的水坝系统化解了暴雨和山洪的危险。云破天青，阳光漫照在平畴绿野、丘陇山林、民居王宫间，城内外又恢复了往日的繁忙。这是一座无与伦比的水上都城，9座城门中有8座均为水门，城内的51条河道和内外护城河绝大多数都是由人工开挖而成。城内外水系相互连通形成了完整的水上交通网络。钟家港古河道也是良渚城这3万多米的人工水道中的一段，在其两岸的部分台地边缘，还发现有保存良好的木构护岸遗迹。河道内的生活废弃垃圾不多，说明当时为保证水上交通进行了较严格的管控。

研究者溯溪而上，探索建城材料的来源。满目青翠间，与良渚人的竹筏隔着5000年的时光错身而过。生长在天目山中海拔四五百米高处的木材，经由水路被运送到良渚城内的钟家港岸边。在这里，麻栎树干被加工成圆木驯，覃树干被加工成方木，它们将被运至莫角山上成为架设神殿的梁柱。

良诸城墙的总长度将近6公里，宽度平均为50米左右。为了增强稳定性，避免地下水对墙体的渗透，人们在纯净黄土夯筑的墙体底部铺垫石块作为基础，这些石料几乎都来自城外西面，东南面的山中。

考古工作者凝视着残余的墙体，分辨着地层，观察着制作技术。城墙之内仿佛站着古时的一群人，他们也曾望向没有竣工的墙体，望向未知的历史那头．.

古城中央那高达10米，面积约30万平方米的莫角山是依托自然体堆铸成的巨型土台。高台之上的三土墩建筑群是大大小小的宫殿，土墩之间7万平方米的夯柱沙土广场上也曾矗立着一些礼仪性的建筑。

良渚的日常居住人口可能是两到三万。莫角山上居住着王和贵族们，宫殿区外围人工堆筑的丘阜上分布着大大小小的作坊，从事玉器、漆器等高端手工业生产的人群居住于此。厚重的城墙形成环抱之势，东北之雉山和西南之凤山各倚一角，状如后世瞭敌警备的角楼。城墙以外自北向南的长条形区域内，数座人工堆筑的土垄上分布着小型聚落。

莫角山南的池中寺遗址出土了相当于 19.7万公斤稻米的碳化稻谷堆积，足够5,000人一季度的消耗。对稻谷的锶同位素分析结果显示，他们来自良渚城以外的不同地方。城内出土的猪骨同样也有多个来源，这表明良渚古城的统治者们可以在大范围内调配基本的食物资源。当原先的城址规模渐趋饱和后，人们扩建了外郭城，历经数百年的开发和建设，良渚城的总面积竟致630万平方米之巨，是那个时代当之无愧的超级"都市"。

良渚古城和水利系统工程的土方量总计约1100万立方米。假使1万人每日风雨无阻的连续施工，也需要11年时间才能完成。

是什么样的力量能聚集这么大范围的人群和资源，建立和维续着如此辉煌的都市？是以信仰为纽带组织起来的权利。

如果我们必须通过空中的上帝视角，才能看清楚良渚人规划的古城和水利之宏大，那么我们需要放大镜，才能观察到良渚信仰的精微与幽深。

玉琮是良渚人至高的通神法器，这件玉丛重6.5公斤，是目前已知的良渚玉琮之首。它出自反山墓地12号墓，墓主应是某代良渚王。玉琮四面的直槽内都以减地微雕而成，一上一下两个神人兽面纹，只有3厘米高，4厘米宽。其间雕琢极为细密，一毫米内就刻画了四五道线条；玉琮的四角则刻有上下分节的神人和兽面；两侧是抽象的侧身飞鸟纹。

这是良渚人心中的神明，他头戴"介"字形羽冠，耸肩、平臂、屈肘叉腰。腰部的位置有一圆目、阔、带獠牙的兽面，双足似爪。这神人兽面像便是良渚的神徽。

神徽中神人兽面的组合表现变化多样，但万变却又不离其宗。它不仅在玉器上被大量表现，也见于其他材质的载体，如象牙器、漆器、陶器等。

早在良渚建成之前，瑶山、北村、官井头等墓地就已出土了大量的玉器，围绕这些玉器的使用，已经建立起了一套严格的规范，体现出明确的尊卑之别，男女之分。早期女性贵族的地位非常突出，墓中可发现一套套完整的玉璜和圆牌饰组合以及偶见但重要的艺术形象。

良渚的信仰和观念也已形成，龙首、人面、神兽面，成为在各类玉器上反复出现的母题。尽管刻划风格时有拙趣，也常常用到镂空工艺，但同良渚古国发展高峰出现的具象的"神徽"显然是同源同意。

刻纹玉器之于良渚人而言，不仅仅是社会权力的表现，更是精神世界的映照。在良渚的社会结构中鲜明的等级差异贯穿生死。

城外的汴家山上葬着平民，文家山上葬着较低等级的贵族，高级贵族先后葬在城内的姜家山和城外的瑶山、汇观山等高台墓地上。莫角山西侧的反山长垄就曾是数代王族的安葬之处。

考古工作者正在绘图记录着反山第20号墓的情景。这是反山墓地中出土玉器种类最齐全、数量最多的一座墓葬。半圆形饰、成组锥形器、三叉形器和冠状器构成了男性头部玉饰的制度化组合。那由300多颗玉管和玉珠所形成的串饰，并非墓主人生前的日常配饰，而是重大仪式中的特殊用器。墓内的4件玉琮、1件玉钺、24件石钺和43件玉璧，均具有象征性的含义。

琮是良渚信仰体系的载体，钺是男性贵族社会权利的物化表现，玉璧则是财富的象征。种种迹象表明这也是一座王墓，年代稍晚于12号墓。

透过专业的目光和笔触，散落的玉管连成珠串，繁复的玉饰各归其位，叠压的玉璧、石铖依着顺序逐层入葬。它让我们越过时间和尘土，同良渚人一起见证那场贵族的葬礼。

与20号墓大致同时的22号墓，其玉器组合却迥然有别。墓主的头颈部放置一组由12件玉管和1件刻纹半圆牌饰组成的串饰。她的胸腹部纵向放着一件玉璜，一组圆形龙纹牌式的组玉佩。这些事物与一同随葬的冠状器、琮式管、环镯、玉璧、玉鱼等玉器，共同彰显着一位女性墓主的权贵身份。

良渚权贵通过一整套玉礼器及其背后的礼仪系统，显示对神权的控制，进而延伸至对世俗王权、军权和财权的占有。

这套权力信仰体系在临平横山、青浦福泉山、吴县草鞋山、武进寺墩和无锡邱承墩等地的高等级墓葬中呈现出惊人的相似性。

在高度一致的信仰系统和文化认同中，各区域中心之间交流分配着象征权力与信仰的高等级玉器和图像。由此建立良渚古国间的政治和宗教联系，构建稳定的社会权利网络，创造了整个良渚文化区域内的共同秩序。毫无疑问，良渚遗址群为代表的社会组织已经进入早期国家的文明阶段。

距今 4,900 年前，深秋黎明的莫角山上，庄严的神殿浸染在霞光中。神殿广场之上已聚集了来自良渚古国各地的首领。一场盛大的祭祀仪式即将开始。

大巫师上前将玉琮交到良渚王手中。良渚王头戴玉锥形器冠，额束缝缀半圆形神人兽面纹玉牌的丝带，发髻上插着三叉形器与长玉管组成的玉饰，脖项间挂着玉管串饰。并肩立于他身旁的王后，髻上插着玉梳背，头上戴着半圆形玉牌与玉管串成的冠饰，颈间挂着玉璜和圆牌饰组成的繁复组佩，双腕戴着玉镯，行动间琳琅珑璁。

良渚王郑重地举起玉琮，那一刻玉琮上的神徽无比清楚而神圣，仿佛神的眼睛正注视着良渚古国的芸芸众生。

通过开创性的政治实践，长江下游的先民缔造了良渚古国。与此同时长江中游的先民们也在探索着区域性整合的道路。

最早铸造城墙的城头山，在距今5,300年前扩建为规整的圆形城池。高耸的城墙，宽深的护城河，这些不再只作为排涝防洪的设施，同时也成为了严密的军事防御屏障。

城头山扩建后不久，在16公里外一座新城出现，这就是鸡叫城。起初居民依照传统修建了环壕聚落，直到距今5000年前随着聚落规模的扩大，人们决定拓展居住空间。他们于是填平环壕，在更外围开挖护城河，挖出的土堆在护城河内侧，便筑起宽阔的城垣。

从普通环壕聚落演变为城池，不仅意味着生活空间的拓展和改造，更意味着通过城壕等大型公共设施的修建，发展了在大规模人群内部进行协作的机制。

随着鸡叫城的发展壮大，周边部落纷纷归附，城外不断出现新的居民点，这些蜂拥而至的人群造成了粮食紧缺的局面。城中的长老们为此争论不休，有人提出应强制驱赶一部分城外的居民点，有人提议抢占青河城和走马岭城的土地，也有人主张扩大耕田面积并将粮食进行集中分配。但旁人均觉得此法见效太慢，于是这声音便被淹没在一片喧扰间。

原野的高阜上，城郭剪影浮现，青年首领走向那位略显落寞的长老。这片肥沃土地是先祖贻留的宝贵遗产，值得为之殚心竭虑，矢志不渝。

不久后，在鸡叫城的周边，人们丈量土地，规划空间，开沟掘土。一条条相互平行的沟渠在大地上铺展开，所经之处便是新造的水田。

考古发掘揭露了城中一片距今4,800年左右的谷糠堆积层，这层谷糠的堆积厚度约15厘米，揭取的面积为80平方米，可大致推算出其代表的稻谷重量约为2.2万公斤。而这仅仅是本地稻作农业规模的冰山一角，长期持续的农田水利建设彻底改变了周围的大地景观。数百年后，鸡叫城及周边区域形成了以城墙为中心的三重环壕。第一重环壕之内是城的主体，护城河和第二重环壕之间为居住和生活区，二重环壕之外是稻田耕作区，一道道平行水渠贯通着环壕之间的水系。

随着社会生产和生活的空前繁荣，鸡叫城得以牢牢控制周边的几十处从属聚落，取代城头山成为整个澧阳平原的中心聚落。

如果说古国是一种文明形态，那鸡叫城就是从自身泥土和稻田里生长出来的农业文明。鸡叫城的社会组织行为开启了新的发展阶段，多组可能用于公共活动的大型建筑遗迹的出现就是明证。

其中63号房址的保存完好度最为令人吃惊，房屋由主体建筑和外围廊道组成。从东到西面阔5间，室内建筑面积420平方米，加上廊道总面积为630平方米左右。先民开挖基槽后，垫长木板以作基础，于木板一侧立柱，有些木板边还可见抬板时留下的绳索。木柱极为考究，以直径约0.5米的半圆形大木柱为主体，间以长方形小木柱，两侧均有斜穿孔。这座距今4,700年左右的史前大宅，间架有序，布局规整，是先民凭着仅有的石制、木制工具以及他们的双手和巧思建造起来的。中国木构建筑传统的特色已然尽显其中。

城头山和鸡叫城虽然是澧阳平原上两个重要的中心聚落，但若从整个长江中游来看，它已远远无法与江汉平原上正迅速崛起的诸城相匹敌。

汉江北岸的大洪山南麓，屈家岭文化的早期城址聚落群在5,300年前便初具规模。直到500年后那始于青萍之末的溪洞终于酿成壮观的风潮。

石家河城扩展成占地达120万平方米的大城，周边的直接控制区域更是足有8平方公里之广。以石家河城为首，江汉之间的聚落进入新一轮的发展高峰。残留于地面上的西城墙曾经在护城河中流淌的水，还能让我们依稀看见四五千年前石家河古城的巍峨。

在石家河城引领下的各个次级城址，如京山屈家岭、应城门板湾、安陆王谷溜、沙洋城河、天门笑城和孝感叶家庙，无不规模可观。

在长江中游一个强大的古国联盟正在形成，要凝聚这样的大型政治实体与社群，必然需要社会意识形态的革新与再造。

石家河城址中央台地上的谭家岭曾发现过大型建筑的遗存，或许是统治者曾经的居处。西北角的邓家湾是一处祭祀场所，在石家河城未起时就存在，城起之后仪式不断。

陶人偶和陶塑动物出于低洼之处和祭祀坑之内，足有上万件，或许象征着献祭之牲。人偶抱着鱼跪坐于地，如虔诚奉献。这样的大型祭祀场所不止一处，城外西侧的印信台上有5个人工堆筑的台基。平整过的地面上密布着祭祀坑，坑内埋藏着套接在一起的大口缸，多的可达25千，延续10米之长。

刻符陶缸起源于遥远的山东地区，大汶口人通常将其用于随葬。石家河在远距离文化交流中并未直接因袭，而是结合本地的实际需求进行了改造，使之成为凝聚本邦各区域政治实体的仪式工具。

站在石家河城内西南隅的三房湾遗址旁，很难不被这密集堆放的红陶杯震惊。它们厚达数米，总数可以百万计。这些陶杯的形态不尽相同，以斜腹杯、斜弧腹杯、喇叭口卷沿杯为主。同一地点分布的黄土坑、烧土面和洗泥池，显示这里曾是专门烧制红陶杯的窑厂，且沿用了数百年之久。这大规模的专业生产线折射出石家河社会的资源掌控能力，工艺制作实力以及举办大型仪式的频繁需求。

4300年前的石家河城，声势浩大的盟会始于一场庄重的祭祀仪式。长江中游诸城的首领们跟随石家河君长步入会场。君长将红陶杯中的酒酶于地面。

巫师将桃人偶和陶塑动物放于预先挖好的祭祀坑前，以敬奉祖先与神灵。首领们郑重捧起刻有本城符号的套缸，与他城的套缸口底相接，象征着参盟者紧密团结，休戚与共。

长江中游的先民们在江汉沅澧间扎根，农作已数千年。由于人口增多，交流频密且河湖交错，洪涝频繁，一个个规模庞大人口集中的古城在湖沼边缘的平原地带拔地而起，呈串联之势。通过仪式性的联盟、公共性的协作，层级城市体系的构建，他们凝聚起共识，认同于彼此。用千百年的时间共同耕耘一片土地，用千万双手共同搭建一座座城，直到缔结一个稳定的区域性政治联盟，造就一方土地的强盛。

盛夏的岱宗之北，古济水之滨，焦家聚落周围，流水潺湲，林木葱郁，座座村舍升起炊烟袅袅。

大概在距今5000年前后的大汶口文化中期，焦家先民开壕筑城，成为目前考古所见海岱地区年代最早的城。有别于良渚古国中贵族与平民异地而葬的制度，焦家遗址目前发现的近400座墓葬，大小墓葬共处于一个墓地，按等级分区，排列相对井然有序。

152号墓葬的墓主是一位老年男性，墓葬占地近12平方米，使用了两椁一棺作为葬具，是目前发现最早的具有三重棺椁的墓葬。内层棺中墓主头部有束发的骨梳和耳坠，手臂佩戴着玉环臂，手指上套有玉指环，下腹部放置象征军权和王权的多孔玉刀、装柄玉钺，身侧和小腿部位放置着作为法器或响器使用的龟甲器。

棺椁之间随葬了27件陶器，有陶鼎、陶豆等炊煮盛器，还有成组的背壶和高柄杯等水具饮器。功能各异的陶器被有序地陈列在独属于死者的空间内。

对于焦家贵族而言，死亡并非终结而是延续，他们希望通过这些沉默的器物展现出繁复讲究的宴飨仪轨，并将那雍雍肃肃的隆重，从生前一直贯穿至地下。

与这些复杂而奢华的大中型墓形成强烈对比的则是那些几近一无所有的小墓。焦家遗址所显露的阶层分化现象意味着大汶口文化中晚期的鲁北地区，已经开启了社会复杂化的进程。由焦家遗址向南约200多公里，鲁中南山地的西南麓是山区至平原的过渡地带。这里背倚青岭，沃野千里，古泗水静静流淌。

距今 5000年前后，岗上城之始建，最终发展成为海岱地区这一时期最大的城，统领周边60余处聚落。

在岗上遗址的大汶口文化晚期墓地中，考古工作者发现了成排排列的墓葬，最为特殊的是一个合葬大墓，面积达10.56平方米，四位墓主均为男性，自北向南年龄递减。葬具是并排的三联棺，带有头箱和边箱。

这座墓随葬的陶器数量达到300多件，以壶、杯、豆等宴饮器为主，而觚和大口尊等礼器则是墓主显要身份的象征。墓中4人均随葬玉钺，多放置于腰胯附近，但这些玉钺在形制、大小、玉料材质等方面判然有别，又体现出同一阶层内不同墓主间也存在明显的等级差异。女性大墓的体量和随葬品与男性大墓相当，但葬具为一棺，头部有骨梳，颈部放着一串绿松石串饰，随葬玉器多为饰品。二层台西南角放置大量猪上颌骨、猪蹄，是墓主生前财富的重要表现形式。

大汶口各区域中心的大墓虽各具特色，但其统一且制度化的一面已经成型，由此开启了中国特有的"礼制"和葬仪。

例如营建大型墓穴和层层相套的多重棺椁，又如以配玉和玉兵等玉器的使用，实现对个人身份的认定与表达。而在棺椁间依次放入成组成套建制规整的陶器组合，则是具有公共性的仪式活动和社会地位的展现。

在后世儒家的礼治中将社会秩序和伦理以物化的形式体现在丧葬仪式活动中，对自天子至庶民各等级死后的棺椁层数，随葬品的种类数量和摆放方式等都有明确的规定。

而 5000年前海岱山岳各处的大汶口高等级墓葬已经拥有这样一套具备严密秩序的丧葬礼仪。在死后长眠处，大汶口文化人群表现出对于统一社会规则和理念的强烈依从性。以死观生，生前的组织严密性更是可想而知。

黄河流域关于古国的探索并未止步于早期礼制的构建。一场更为深刻的社会变革即将掀开新章，被后来的考古学家们称之为"龙山"的崭新时代喷薄欲出。

4500年前一个盛夏的晌午，日照海滨的阳光灼热而刺眼。丹土城的大殿中贵族之间正发生一场激烈的争吵，穿越大汶口文化晚期此时的丹土已有百年建城史。社会分层固化，年长的贵族们依靠血缘关系结成的权力网络大大的限制了酋长的权力，也压制了反对者的声音。与丹土城毗邻的尧王城，无论在选址还是规模上都不可小觑。团结奋进的气势更是让丹土年轻人羡慕，或许正是鲁东南地区这些早期复杂社会的发展与竞争，促进了大汶口文化迈向龙山时代的社会变革。

要获得更为充分的社会掌控权和财产权，自龙山时代开始的分家在所难免。酋长终于决定带领着年轻人去往丹土以东4公里外的两城镇，踌躇满志地要在这黄海之滨开辟一片新的天地。山东龙山文化从大汶口文化脱胎出来的过程中，聚落数量快速增加，在一些地区呈现几何量级的增长，聚落规模也在快速扩张。

鲁北的寿光边线王、临淄桐林、邹平丁公、章丘城子崖和鲁西的景阳冈，鲁东南的五莲丹土，日照两城镇。尧王城和鲁南苏北的藤花落都是所在区域的核心城址。串起这两个城市群的应该是当时的主要交通路线。

龙山先民不再像大汶口时代那般，营建以血缘原则安排的大规模族墓地，这或许又意味着城内的居民们已处于杂居状态，最初的城市生活开始了。

黄海之滨的两城镇在4,000多年前的改革者的推动下，成为现在所能见到面积最大的龙山文化遗址之一。大汶口人在手工业经济领域的优势已被龙山县民继承，并发展出初级专业化和社会化的标准。

两城镇便是一处以制造业见长的中心聚落，作为当时最发达的工业产品，陶器有着巨大的产量和复杂的器类，其中蛋壳陶就代表了当时最顶尖的制陶技术。

色墨如漆，光泽莹润，器体轻薄，质硬如瓷这便是龙山文化的薄胎黑陶杯，因壁薄如蛋壳而得名"蛋壳陶"。一件这样的杯子，最薄处的杯口仅有0.2毫米，重量不过93克，即便以现代工艺也难以完全复制。杯柄中部鼓出部位中空并装饰细密的镂孔，其内还放置一粒陶丸，将杯子拿在手中晃动时，陶丸碰撞笼壁会发出清脆的响声。它的造型挺拔秀丽，脆弱又精致，如用墨色晕染的一株含苞欲放的百合，静静伫立在时光的薄暮中……

核心城址依赖达的制陶业等手工业经济成为区域的中心。周边的小型城市和聚落对其产品存在需求，于是产生了频繁的交换，乃至商贸活动。在对这些经济活动进行管理时，便诞生了制度与权利。

随着城市社会的产生，社会分化进一步加剧。玉器是当时的高端产品，目前发现的大部分玉器，比如多孔刀、钺、璇玑、璧、琮、璋、圭等礼器，基本都出于大型的中心聚落或是祭祀遗迹中，应该属于专供上层或仪式使用的奢侈品。

而位于山东省中部地区的临胸西珠峰遗址发现了目前考古所见，龙山时代规格最高的大墓。墓内所处的两只玉笄精美绝伦，那种处于整个社会最顶端，可能具有"王"的身份的阶层已经产生。

在邹平丁公遗址出土的一片大平底盆底部残片上，俨然可见线条流畅，排列规律的五行 11个符号，刻写有一定章法，为探讨原始文字的产生提供了珍贵的实物资料。

古国时代列城林立的景观正与先秦文献中上古时代天下万国的记忆合节。

东部地区良渚古国以信仰为基石，以权力为利器一统长江下游，已然迈入文明的门槛，成为东亚最早的区域性国家。

在良渚人以神权维系长江下游古国时，黄河下游的海岱地区却走向了更为世俗化的社会。他们自现实生活中发展出社会治理的智慧，较之鬼神却更崇人力，继承并发展更为严密的礼制实践，成为了"礼出东方"的强大力量。

石家河扎根于传统，在长江中游建立起强大的连城网络，形成千里同风的联盟。大约在4,100年前来自北方的力量南下，改变了石家河文化的面貌，长江中游也由此纳入早期中国的版图。后石家河文化中的玉器和图像元素均在日后的跨区域交流中扮演重要的角色。

除此之外所有地区都在不同程度地参与着此后作为整体的华夏文明的构建。正是他们的发展造就了广域中国的多元基础，又为王朝中国的延绵发展提供了生态、文化、人口和资源等各方面腾挪的空间。

黄河之水奔腾而下，中原即将成为新的时代引领者，下一场历史变革就要到来。